% !TEX program = lualatex
%================================================================
% 中間発表 計画書・設計書 — チーム23
%
% 編集規約・工程表: report/計画書_中間発表/_作業計画.md
%   - 図は fig/*.drawio を draw.io CLI で PNG 化し \includegraphics で挿入
%   - 完成目標: PDF 50 ページ以内（本文 47 ページ目標）
%   - 各 \section 直下の「TODO(Phase N)」は作業計画 §6・§7 への対応
%================================================================
\documentclass[lualatex,book,paper=a4paper,jafontsize=11.5pt,jafontscale=1.05]{jlreq}

% --- フォント・表紙（hack-*.sty を流用）---
\usepackage{luatexja-fontspec}
\usepackage{hack-fonts}
\usepackage{hack-cover}

% --- 図・表 ---
\usepackage{graphicx}
\graphicspath{{fig/}}
\usepackage{float}
\usepackage{array}
\usepackage{tabularx}
\usepackage{booktabs}
\usepackage{multirow}
\usepackage{longtable}

% --- コード・整形 ---
\usepackage{listings}
\usepackage{fancyvrb}
\usepackage[most,breakable]{tcolorbox}
\usepackage{amsmath}
\usepackage{enumitem}
\usepackage{url}
\usepackage{needspace}
\usepackage{placeins}
\usepackage{xcolor}

% --- 作図（TikZ：正常な既存図は据え置くため残す）---
\usepackage{tikz}
\usetikzlibrary{arrows.meta,positioning,shapes.geometric,fit,calc}

% --- リンク ---
\usepackage[hidelinks]{hyperref}

% --- ユーザ定義 ---
\newcommand{\code}[1]{\texttt{#1}}
\definecolor{lightgray}{HTML}{F2F3F4}
\definecolor{codebg}{HTML}{F7F9FB}

\lstdefinestyle{textdiagram}{
  basicstyle=\ttfamily\small, numbers=none, breaklines=true,
  frame=single, tabsize=2, showstringspaces=false
}
\lstdefinestyle{pseudo}{
  basicstyle=\ttfamily\small, numbers=left, numberstyle=\tiny,
  breaklines=true, frame=single, tabsize=2, showstringspaces=false
}

% 改ページ制御（孤立行・寡婦行の抑制）
\widowpenalty=10000
\clubpenalty=10000
\predisplaypenalty=10000

% --- 表紙メタ情報 ---
\title{ハッカソン}
\subtitle{～計画書と設計書～}
\team{チーム23}
\author{
  24G1038 : 片岡 聖\\
  24G1175 : 地曵 賢人\\
  24G1131 : 御代川 稜\\
  25G1021 : 梅澤 颯太\\
  25G1053 : 齋藤 翔太\\
  25G1065 : 塩澤 匠生
}
\date{\today}
\version{1.0}

%================================================================
\begin{document}
\maketitle

\setcounter{tocdepth}{2}
\tableofcontents \thispagestyle{empty}

%================================================================
% Chapter 1 — 計画書           ［Phase 1］
%   入力: 23_ 行 139–500 / docs team/schedule.md・roles.md /
%         review_プロジェクト見直_片岡.md
%================================================================
\chapter{プロジェクトの計画書}\label{chap:project-plan}\setcounter{page}{1}

\section{プロジェクトの概要}\label{sec:overview}
\subsection{プロジェクトの題目}
% TODO(Phase 1): 議事録 9–10 章準拠。散文で簡潔に（予算 2 頁）
\subsection{プロジェクトの目的}
% TODO(Phase 1): 解決すべき課題と最終ゴール
\subsection{既存の技術とプロジェクトの独自性}
% TODO(Phase 1): 既存技術との差分

\section{スコープ・マネジメント}\label{sec:scope}
\subsection{成果物}
% TODO(Phase 1): 成果物一覧（予算 4 頁：本節〜FBS/PBS 対応まで）
\subsection{プロジェクトの対象範囲と非対象範囲}
\subsection{機能分解構造FBS}
% TODO(Phase 1): fig/fbs.png（draw.io 再作成）
\subsection{製品分解構造PBS}
% TODO(Phase 1): fig/pbs.png（draw.io 再作成）
\subsection{FBSとPBSの対応関係}
% TODO(Phase 1): 小さな対応表 1 つ

\section{作業計画}\label{sec:workplan}
\subsection{作業分解構造WBSと管理番号}
% TODO(Phase 1): fig/wbs.png（draw.io 新規。表で足りれば表化可。予算 5 頁）
\subsection{アローダイアグラムとクリティカルパス}
% TODO(Phase 1): fig/arrow.png（draw.io 新規。クリティカルパス強調）
\subsection{役割分担}
% TODO(Phase 1): 役割分担表 1 つ
\subsection{ガントチャート}
% TODO(Phase 1): fig/gantt.png（draw.io 新規。4 フェーズ・WBS 番号対応）

\section{検証と妥当性確認: V\&V計画}\label{sec:vv}
% TODO(Phase 1): 予算 2 頁・表中心。MOP 階層は 1 表に集約
\subsection{プロジェクトの成果物に対する有効指標MOE}
\subsection{有効指標MOEに対応する性能指標MOP}
\subsection{システムテストで評価する性能指標MOP}
\subsection{結合テストで評価する性能指標MOP}
\subsection{単体テストで評価する性能指標MOP}
\subsection{プロジェクト中に監視する技術性能指標TPM}
\subsection{プロジェクト中に監視する指標TPM}

\section{資源・調達管理}\label{sec:resource}
% TODO(Phase 1): 1 頁目安

\section{リスク管理}\label{sec:risk}
% TODO(Phase 1): 1 頁目安

%================================================================
% Chapter 2 — システムの基本設計   ［Phase 2：Block A/B 重複解消］
%   入力: 23_ 行 501–1209 / .agent/architecture.md / review_arduino_*.md
%   方針: Block B（\label 付き）を土台に Block A 固有事実を畳み込む
%================================================================
\chapter{システムの基本設計}\label{chap:basic-design}

\section{機能一覧}\label{sec:basic-functions}
% TODO(Phase 2): 23_ 行 506(A)/817(B) を 1 版へ。要求一覧＋FBS（予算 3 頁）

\section{システム構成図}\label{sec:basic-architecture}
% TODO(Phase 2): 23_ 行 569(A)/881(B)。全体ブロック図 1 枚＋短い説明（予算 3 頁）

\section{操作インターフェース設計}\label{sec:basic-ui}
% TODO(Phase 2): 23_ 行 672(A)/993(B)。LED 表示・ユーザー操作を各 1 表（予算 2 頁）

\section{状態遷移図}\label{sec:basic-state}
% TODO(Phase 2): 23_ 行 721(A)/1122(B)。指揮者・楽器の状態遷移図＋定義表（予算 4 頁）

%================================================================
% Chapter 3 — システムの詳細設計   ［Phase 3：重複解消＋最大の削減対象］
%   入力: 23_ 行 1210–3056 / .agent/api.md / architecture.md
%   方針: 長い API 表・コード例は要点の散文＋小図表に置換
%================================================================
\chapter{システムの詳細設計}\label{chap:detail-design}

\section{部品一覧}\label{sec:parts-list}
% TODO(Phase 3): 23_ 行 1214(A)/1873(B)。PBS 要素を 1–2 表（予算 2 頁）

\section{ハードウェア設計}\label{sec:hw-design}
% TODO(Phase 3): 23_ 行 1219(A)/1932(B)。配線・ピン表＋接続図 1 枚（予算 3 頁）

\section{ソフトウェア設計}\label{sec:sw-design}
% TODO(Phase 3): 23_ 行 1308(A)/1997(B)。クラス図＋パケット定義＋楽譜形式（予算 6 頁）
\subsection{ファイル構成}
\subsection{オブジェクト設計(クラス図)}
\subsection{関数設計}

\section{処理フローの設計}\label{sec:proc-flow}
% TODO(Phase 3): 23_ 行 1716(A)/2467(B)。3 フェーズモデル＋処理フロー図（予算 3 頁）

\section{テストの設計}\label{sec:test-design}
% TODO(Phase 3): 23_ 行 1834(A)/2959(B)。設計目標値表＋検証方針（予算 2 頁）

%================================================================
% Chapter 4 — 生成 AI の利用        ［Phase 4］
%   入力: 23_ 行 3058–3077（流用・微圧縮）
%================================================================
\chapter{生成AIの利用}\label{chap:genai}
% TODO(Phase 4): 既存記述を流用・微圧縮（予算 1 頁）

\end{document}
